% ==========================================
% CHAPTER 1: VISCOSITY & PROBLEM SOLVING
% ==========================================
\section{Viscosity and Shear Stress: Problem Solving}

\subsection{Problem 1: Plate moving between two fluids}

\subsubsection{Problem Statement}
A thin flat plate is moving horizontally at a constant, unknown velocity $V$ between two other parallel plates. The bottom plate is stationary, while the top plate moves at a constant velocity $U$. 
The space is filled with two different Newtonian fluids:
\begin{itemize}
    \item \textbf{Fluid 1 (Bottom):} Dynamic viscosity $\mu_1$, filling a gap of height $c \cdot t$.
    \item \textbf{Fluid 2 (Top):} Dynamic viscosity $\mu_2 = 2\mu_1$, filling a gap of height $(1-c) \cdot t$.
\end{itemize}
\textbf{Goal:} Determine the steady-state velocity vector $\vec{V}$ of the middle plate.

\begin{center}
\begin{tikzpicture}[scale=1.2, >=Stealth]
    % Total height t = 4 for visual purposes. Let's say c = 0.35. h = 1.4
    \def\totH{4}
    \def\midH{1.4}
    
    % Bottom Plate (Stationary)
    \fill[pattern=north east lines, pattern color=gray] (0,-0.3) rectangle (7,0);
    \draw[thick] (0,0) -- (7,0) node[right, font=\bfseries] {$v = 0$};
    
    % Top Plate (Moving at U)
    \draw[thick, primary, ->] (2,\totH) -- (6,\totH) node[right, font=\bfseries] {$U$};
    \draw[thick] (0,\totH) -- (7,\totH);
    
    % Middle Plate (Moving at V)
    \draw[thick, accent, ->] (2,\midH) -- (5,\midH) node[right, font=\bfseries] {$V$};
    \draw[thick, accent] (0,\midH) -- (7,\midH);
    
    % Fluids
    \node at (3.5, 0.7) {Fluid 1: $\mu_1$};
    \node at (3.5, 2.7) {Fluid 2: $\mu_2 = 2\mu_1$};
    
    % Dimensions
    \draw[<->] (-0.5, 0) -- (-0.5, \midH) node[midway, left] {$c \cdot t$};
    \draw[<->] (-0.5, \midH) -- (-0.5, \totH) node[midway, left] {$(1-c) \cdot t$};
    \draw[<->] (-1.5, 0) -- (-1.5, \totH) node[midway, left] {$t$};
    
    % Velocity Profiles (Linear)
    \draw[dashed, gray] (0,0) -- (2,\midH); % Bottom profile
    \draw[dashed, gray] (2,\midH) -- (4,\totH); % Top profile
    
    \foreach \y in {0.2, 0.4, 0.6, 0.8, 1.0, 1.2, 1.4} {
        \draw[->, gray, very thin] (0,\y) -- ({\y*1.428},\y);
    }
    \foreach \y in {1.6, 1.8, 2.0, 2.2, 2.4, 2.6, 2.8, 3.0, 3.2, 3.4, 3.6, 3.8, 4.0} {
        \draw[->, gray, very thin] (0,\y) -- ({2 + (\y-1.4)*0.769},\y);
    }
\end{tikzpicture}
\end{center}

\begin{notebox}[Scientific Verification: Linear Velocity Profile]
The assumption that the velocity profile $u(y)$ is strictly linear between the plates is not arbitrary. It is mathematically derived from the \textbf{Navier-Stokes equations} (formulated by Navier in 1822 and Stokes in 1845). For a steady, fully developed, one-dimensional laminar flow with zero pressure gradient ($\frac{\partial p}{\partial x} = 0$), the equation simplifies to $\frac{d^2u}{dy^2} = 0$, whose integration strictly yields a linear profile $u(y) = Ay + B$. This specific setup is historically known as \textbf{Couette Flow}, experimentally validated by \textbf{Maurice Couette in 1890}.
\end{notebox}

\subsubsection{Governing Equations \& Force Balance}

Consider a small piece of the middle plate of area $A$. The plate experiences a shear stress from Fluid 1 on its bottom surface ($\tau_1$) and a shear stress from Fluid 2 on its top surface ($\tau_2$). 

Since the plate moves at a \textit{constant velocity} $V$, its acceleration is zero. According to Newton's First Law of Motion (1687):
\begin{align*}
    \Sigma F_x &= m \cdot a_x = 0 \\
    F_{\text{shear, top}} - F_{\text{shear, bottom}} &= 0 \\
    \tau_2 \cdot A - \tau_1 \cdot A &= 0 \implies \tau_1 = \tau_2
\end{align*}
\textit{(Note: Original handwritten notes contained a transcription error stating $\frac{\tau_2}{t} - \frac{\tau_1}{A} = 6$. This is mathematically and dimensionally incorrect. The derivation above represents the correct physical principle.)}

Using Newton's Law of Viscosity ($\tau = \mu \frac{du}{dy}$), we determine the shear stresses. Because the velocity profiles are linear, the velocity gradients $\frac{du}{dy}$ are constant in each layer:

\begin{align}
    \tau_1 &= \mu_1 \frac{\Delta u_1}{\Delta y_1} = \mu_1 \frac{V - 0}{c \cdot t} = \mu_1 \frac{V}{ct} \\
    \tau_2 &= \mu_2 \frac{\Delta u_2}{\Delta y_2} = \mu_2 \frac{U - V}{(1-c) \cdot t} 
\end{align}

\subsubsection{Algebraic Resolution}

Substitute $\mu_2 = 2\mu_1$ and equate the two shear stresses ($\tau_1 = \tau_2$):
\begin{equation*}
    \mu_1 \frac{V}{ct} = 2\mu_1 \frac{U - V}{(1-c)t}
\end{equation*}

We can cancel out $\mu_1$ and $t$ from both sides:
\begin{equation*}
    \frac{V}{c} = \frac{2(U - V)}{1-c}
\end{equation*}

Cross-multiply to solve for $V$:
\begin{align*}
    V(1 - c) &= 2c(U - V) \\
    V - Vc &= 2Uc - 2Vc \\
    V - Vc + 2Vc &= 2Uc \\
    V(1 + c) &= 2Uc
\end{align*}

\begin{eqbox}{Final Analytical Solution}
\begin{equation}
    V = \frac{2Uc}{1+c}
\end{equation}
\end{eqbox}

\subsubsection{Verification of Physical Limits}

To ensure our formula is robust, we must evaluate its behavior at the geometric boundaries ($c \to 0$ and $c \to 1$). 
\textit{(Note: The original notes contained calculation errors regarding these limits, stating $V=3$ and $V=0$. The scientifically verified limits are detailed below).}

\begin{itemize}
    \item \textbf{Limit 1: $c \to 0$} \\
    If $c$ approaches $0$, the middle plate is effectively pressed against the bottom stationary wall. 
    \begin{equation*}
        \lim_{c \to 0} V = \frac{2U(0)}{1+0} = 0
    \end{equation*}
    This strictly satisfies the \textbf{No-Slip boundary condition} (Stokes, 1851): the velocity of the plate must match the velocity of the bottom wall ($0$).

    \item \textbf{Limit 2: $c \to 1$} \\
    If $c$ approaches $1$, the middle plate is effectively pressed against the top wall which is moving at velocity $U$.
    \begin{equation*}
        \lim_{c \to 1} V = \frac{2U(1)}{1+1} = \frac{2U}{2} = U
    \end{equation*}
    Again, this perfectly satisfies the No-Slip boundary condition relative to the top moving wall.
\end{itemize}